\heading{实验物理中的统计方法-模拟题2-期中}
\noteinfo{题目和答案由AI生成。知识点范围按现有往年题整理，叙述风格尽量向往年题靠拢；本套难度较往年期中略高。}

\begin{question}
$X\sim N(0,1)$，$Y\sim N(0,1)$，并且
\[
\V(2X-Y)=1.
\]
求 $(X,Y)$ 的协方差矩阵。
\begin{solution}
由
\[
\V(2X-Y)=4\V(X)+\V(Y)-4\cov(X,Y),
\]
代入 $\V(X)=\V(Y)=1$ 得
\[
1=4+1-4\cov(X,Y).
\]
故
\[
\cov(X,Y)=1.
\]
因此
\ansmath{\Sigma=
\begin{pmatrix}
1&1\\
1&1
\end{pmatrix}}
\end{solution}
\end{question}

\begin{question}
一个显示屏每次等概率显示 $1$ 到 $N$ 中的一个整数。某人连续看了 $n$ 次，记录下看到的最大值 $K$。求 $K$ 的概率分布，并写出 $E[K]$ 的求法。
\begin{solution}
先有
\[
P(K\le k)=\left(\frac{k}{N}\right)^n,\qquad k=1,2,\cdots,N.
\]
因此
\[
P(K=k)=P(K\le k)-P(K\le k-1)
=\frac{k^n-(k-1)^n}{N^n},\qquad k=1,2,\cdots,N.
\]
期望可由尾和公式写成
\[
E[K]=\sum_{k=1}^{N}P(K\ge k)
=\sum_{k=1}^{N}\left[1-\left(\frac{k-1}{N}\right)^n\right].
\]
也可写作
\ansmath{E[K]=N-\frac{1}{N^n}\sum_{j=0}^{N-1}j^n}
\end{solution}
\end{question}

\begin{question}
一种原子的半衰期为 $8$ 年。现有 $5$ 个原子，经过 $24$ 年后恰好有 $2$ 个尚未衰变的概率是多少？
\begin{solution}
单个原子经过 $24$ 年仍未衰变的概率为
\[
p=\left(\frac12\right)^3=\frac18.
\]
故尚未衰变的原子数服从参数为 $(5,\frac18)$ 的二项分布。于是
\[
P=\binom52\left(\frac18\right)^2\left(\frac78\right)^3.
\]
即
\ansmath{P=\binom52\left(\frac18\right)^2\left(\frac78\right)^3=\frac{1715}{16384}}
\end{solution}
\end{question}

\begin{question}
设 $U\sim U(0,1)$，定义
\[
Y=\tan\!\left(\pi\left(U-\frac12\right)\right).
\]
求 $Y$ 的概率密度函数，并说明 $E[Y]$ 是否存在。
\begin{solution}
令
\[
V=\pi\left(U-\frac12\right),
\]
则
\[
V\sim U\!\left(-\frac{\pi}{2},\frac{\pi}{2}\right),\qquad Y=\tan V.
\]
由变量变换，
\[
f_Y(y)=\frac{1}{\pi}\cdot \frac{1}{1+y^2},\qquad y\in \mathbb R.
\]
这就是标准 Cauchy 分布。由于
\[
\int_{-\infty}^{\infty}|y|f_Y(y)\,dy
\]
发散，因此
\anstext{$E[Y]$ 不存在。}
\end{solution}
\end{question}

\begin{question}
一个矩形的边长分别为 $a$ 和 $b$。在矩形边界上独立、均匀地任取两点，求这两点距离平方的期望值。
\begin{solution}
设边界上随机点坐标为 $(X,Y)$，则两点距离平方为
\[
D^2=(X_1-X_2)^2+(Y_1-Y_2)^2.
\]
故
\[
E[D^2]=2\V(X)+2\V(Y).
\]
对边界均匀分布，直接计算可得
\[
E[X]=\frac a2,\qquad E[Y]=\frac b2,
\]
\[
E[X^2]=\frac{a(2a+b)}{6(a+b)},\qquad
E[Y^2]=\frac{b(a+2b)}{6(a+b)}.
\]
代回化简得
\[
E[D^2]=\frac{(a+b)^2}{6}.
\]
因此
\ansmath{E[D^2]=\frac{(a+b)^2}{6}}
\end{solution}
\end{question}

\begin{question}
某仪器在时间窗口 $[0,T]$ 内会随机出现噪声。已知在该时间窗口内恰好出现两个噪声；若这两个噪声的时间间隔小于 $t$（其中 $0<t<T$），就会被误判成一个信号。求误判概率。
\begin{solution}
条件于“恰好两个噪声”后，两个到达时刻可视为两个独立均匀点
\[
U,V\sim U(0,T).
\]
所求即
\[
P(|U-V|<t).
\]
在边长为 $T$ 的正方形区域中，满足 $|u-v|<t$ 的面积占比为
\[
1-\frac{(T-t)^2}{T^2}.
\]
所以
\[
P(|U-V|<t)=1-\frac{(T-t)^2}{T^2}
=\frac{2t}{T}-\left(\frac{t}{T}\right)^2.
\]
\end{solution}
\end{question}

\begin{question}
随机变量 $(X,Y)$ 的联合密度为
\[
f(x,y)=\begin{cases}
cx,&0<y<x<1,\\
0,&\text{其他}.
\end{cases}
\]
\begin{enumerate}
\item 求常数 $c$；
\item 求 $P\!\left(Y<\dfrac{X}{2}\right)$。
\end{enumerate}
\begin{solution}
由归一化条件，
\[
1=c\int_0^1\int_0^x x\,dy\,dx
=c\int_0^1 x^2\,dx
=\frac{c}{3}.
\]
故
\ansmath{c=3}
再算
\[
P\!\left(Y<\frac X2\right)=\int_0^1\int_0^{x/2}3x\,dy\,dx
=\int_0^1 \frac{3x^2}{2}\,dx
=\frac12.
\]
\end{solution}
\end{question}

\begin{question}
某粒子束中电子所占比例为 $10^{-5}$，其余全为光子。粒子通过一个双层探测器时，给出 0、1、2 个信号的条件概率如下：
\[
P(0\mid e)=0.005,\qquad P(1\mid e)=0.015,\qquad P(2\mid e)=0.98,
\]
\[
P(0\mid \gamma)=0.99898,\qquad P(1\mid \gamma)=0.001,\qquad P(2\mid \gamma)=2\times10^{-5}.
\]
\begin{enumerate}
\item 若只看到 1 个信号，该粒子为电子的概率是多少？
\item 若看到 2 个信号，该粒子为电子的概率是多少？
\end{enumerate}
\begin{solution}
记先验
\[
P(e)=10^{-5},\qquad P(\gamma)=1-10^{-5}.
\]
由 Bayes 公式，
\[
P(e\mid 1)=\frac{P(1\mid e)P(e)}{P(1\mid e)P(e)+P(1\mid \gamma)P(\gamma)},
\]
\[
P(e\mid 2)=\frac{P(2\mid e)P(e)}{P(2\mid e)P(e)+P(2\mid \gamma)P(\gamma)}.
\]
代入数据得
\ansmath{P(e\mid 1)\approx 1.50\times 10^{-4},\qquad P(e\mid 2)\approx 0.3289}
\end{solution}
\end{question}

\begin{question}
蒲丰投针实验中，若针长为 $l$，平行线间距为 $a$（且 $l<a$），则
\[
\hat\pi=\frac{2lN}{an},
\]
其中 $N$ 为总投掷次数，$n$ 为针与线相交次数。某次实验中
\[
l=2.5\text{ cm},\qquad a=3\text{ cm},\qquad N=4000,\qquad n=2122.
\]
忽略 $l,a$ 的误差，求 $\pi$ 的估计值及其不确定度。
\begin{solution}
先有
\[
\hat\pi=\frac{2lN}{an}
=\frac{2\times 2.5\times 4000}{3\times 2122}
\approx 3.14169.
\]
把 $n$ 视为二项分布，取
\[
p\approx \frac{n}{N},
\qquad
\sigma_n\approx \sqrt{Np(1-p)}
=\sqrt{n\left(1-\frac{n}{N}\right)}.
\]
由误差传播，
\[
\sigma_{\pi}\approx \left|\frac{\partial \hat\pi}{\partial n}\right|\sigma_n
=\hat\pi\frac{\sigma_n}{n}.
\]
代入可得
\[
\sigma_\pi\approx 0.0467.
\]
\end{solution}
\end{question}

\begin{question}
设 $X\sim \mathrm{Poisson}(\nu_1)$，$Y\sim \mathrm{Poisson}(\nu_2)$，且 $X,Y$ 相互独立。求条件分布
\[
X\mid (X+Y=n).
\]
\begin{solution}
对 $k=0,1,\cdots,n$，
\[
P(X=k\mid X+Y=n)
=\frac{P(X=k,Y=n-k)}{P(X+Y=n)}.
\]
代入泊松分布并化简：
\[
P(X=k\mid X+Y=n)
=\frac{\binom{n}{k}\nu_1^k\nu_2^{\,n-k}}{(\nu_1+\nu_2)^n}
=\binom{n}{k}\left(\frac{\nu_1}{\nu_1+\nu_2}\right)^k
\left(\frac{\nu_2}{\nu_1+\nu_2}\right)^{n-k}.
\]
\end{solution}
\end{question}

\begin{question}
设单个元件的寿命分布函数为 $F(t)$。现有两个相互独立的元件并联，只要有一个元件仍正常，系统就正常工作。求系统寿命 $T$ 的分布函数。
\begin{solution}
并联系统失效，等价于两个元件都已失效，因此
\[
P(T\le t)=P(T_1\le t,T_2\le t)=F(t)^2.
\]
\end{solution}
\end{question}

\begin{question}
设 $X_1,\cdots,X_n$ 独立同分布于 $U(0,\theta)$，记
\[
Y=\max(X_1,\cdots,X_n).
\]
求 $Y$ 的概率密度函数、期望和方差。
\begin{solution}
对 $0<y<\theta$，
\[
P(Y\le y)=\prod_{i=1}^n P(X_i\le y)=\left(\frac{y}{\theta}\right)^n.
\]
故
\[
f_Y(y)=\frac{d}{dy}\left(\frac{y}{\theta}\right)^n
=\frac{n y^{n-1}}{\theta^n},\qquad 0<y<\theta.
\]
进一步有
\[
E[Y]=\int_0^\theta y\frac{n y^{n-1}}{\theta^n}\,dy
=\frac{n}{n+1}\theta,
\]
\[
E[Y^2]=\int_0^\theta y^2\frac{n y^{n-1}}{\theta^n}\,dy
=\frac{n}{n+2}\theta^2.
\]
所以
\[
\V(Y)=E[Y^2]-E[Y]^2
=\frac{n\theta^2}{(n+1)^2(n+2)}.
\]
\end{solution}
\end{question}

